NDNSF provides a general framework for building service-oriented distributed applications, supporting diverse service types and adapting to varied requirements. It enables fine-grained control over which providers can offer a service and which users can invoke it. Through protected Request--ACK exchanges and user-side selection, a user can invoke the same service across multiple providers. In the single-provider wired loss experiment, NDNSF achieves higher completion rates than NSC and rates comparable to gRPC. In the intermittent-availability comparisons, it achieves higher completion rates than fixed-provider configurations. A controlled ablation shows that task-specific feedback improves deadline completion under changing provider capacity. Against baselines with the same preconfigured candidates and the ability to retry another provider, NDNSF shows no consistent completion advantage and has higher invocation latency and traffic. The costs of obtaining and updating candidate lists are outside that comparison. Under one ABE authority, NAC-ABE aggregates an entity's service permissions in one decryption key. At larger permission counts, this reduces initial distribution and key-rotation costs relative to separate ABE keys for each permission. In the evaluated scheme, granting a permission updates one entity's decryption key, whereas revoking all of that entity's service permissions requires key rotation and redistribution to the remaining authorized entities.

Future work will extend NDNSF's APIs to coordinate multiple providers in collaboratively executing complex tasks. We will further develop and evaluate UAV and distributed AI inference applications built on NDNSF.
